\begin{table}[H]
\raggedright
\caption{Baseline characteristics of the analytic sample ($n=31$ with POAM-P data).}
\label{tab:sample}
\footnotesize
\begingroup
\setlength{\tabcolsep}{4pt}
\renewcommand{\arraystretch}{1.08}
\noindent
\begin{tabular*}{\textwidth}{@{\extracolsep{\fill}}lccc@{}}
\toprule
Variable & Mean (SD) & Range & n \\
\midrule
Age (years) & 50.3 (9.0) & 29-74 & 31 \\
BMI (kg/m2) & 27.3 (6.0) & 18-43 & 31 \\
POAM-P Avoidance (0-40) & 22.0 (6.8) & 10-38 & 31 \\
POAM-P Overdoing (0-40) & 25.3 (5.8) & 12-35 & 31 \\
POAM-P Pacing (0-40) & 24.6 (8.2) & 10-38 & 31 \\
Kinesiophobia TSK (17-68) & 36.4 (8.5) & 23-55 & 31 \\
PROMIS pain interference (sum) & 27.6 (6.2) & 16-40 & 31 \\
Self-efficacy (baseline, 1-5) & 3.3 (0.7) & 2-4 & 31 \\
Outcome expectancy (1-5) & 3.8 (0.4) & 3-5 & 31 \\
Action planning (1-5) & 3.1 (1.0) & 1-5 & 31 \\
\bottomrule
\end{tabular*}
\par\vspace{3pt}\noindent\parbox{\textwidth}{\footnotesize\textit{Note.} POAM-P subscales range 0-40. All participants were women. BMI = body mass index; TSK = Tampa Scale for Kinesiophobia; PROMIS = pain interference.}
\endgroup
\end{table}
